\section{Histórico e Desenvolvimento da Computação Autônoma}
\label{sec:historico}

A necessidade de sistemas auto-gerenciáveis não é um fenômeno exclusivo da computação moderna. Um paralelo histórico instrutivo pode ser traçado com a indústria telefônica da década de 1920. Naquele período, o crescimento exponencial no uso do telefone exigia um número cada vez maior de operadores humanos para gerenciar as centrais de comutação manuais. As preocupações com a possível escassez de operadores qualificados para atender à demanda eram sérias. A solução foi o desenvolvimento das centrais automáticas de comutação, que eliminaram a necessidade de intervenção humana nessa tarefa \cite{W1999392360}. A computação autônoma representa, em essência, uma resposta análoga ao mesmo problema estrutural, agora na escala e complexidade dos sistemas de informação contemporâneos.

O termo \textit{autonomic computing} foi criado pela IBM em 2001. A palavra ``autônomo'' tem origem na biologia: no corpo humano, o sistema nervoso autônomo regula funções inconscientes como o tamanho da pupila, as funções digestivas, a frequência respiratória e a dilatação dos vasos sanguíneos, sem que a mente consciente precise intervir \cite{W1999392360}. A computação autônoma busca intervir nos sistemas computacionais de maneira análoga, reduzindo o envolvimento humano na administração cotidiana. Embora o termo tenha sido formalizado pela IBM, os conceitos subjacentes ao auto-gerenciamento já vinham sendo explorados em diferentes contextos ao longo dos anos anteriores.

\subsection{Projetos Precursores}
\label{subsec:projetos-precursores}

O campo da computação autônoma foi moldado por uma série de projetos independentes que, sem coordenação explícita, convergiram para os mesmos princípios de auto-gestão \cite{W1999392360}.

O primeiro projeto notável foi o \textbf{Situational Awareness System (SAS)}, iniciado em 1997 pela agência de defesa norte-americana DARPA, como parte do programa \textit{Small Units Operations} (SUO). Seu objetivo era criar dispositivos pessoais de comunicação e localização para soldados em campo de batalha. O sistema deveria propagar autonomamente informações de status entre os dispositivos, adaptar-se às mudanças de topologia da rede e ajustar frequências e largura de banda de acordo com as condições ambientais --- inclusive com equipamentos de interferência inimiga em operação. O projeto pioneirou o roteamento ad-hoc descentralizado e adaptativo como forma de auto-gerenciamento em ambientes adversos, lidando com redes que podiam chegar a 10.000 nós móveis \cite{W1999392360}.

Em 2000, a DARPA lançou o programa \textbf{DASADA} (\textit{Dynamic Assembly for Systems Adaptability, Dependability, and Assurance}), cujo objetivo era desenvolver tecnologia para sistemas críticos com alta confiabilidade e adaptabilidade. Este projeto introduziu a abordagem orientada a arquitetura para o auto-gerenciamento, com a noção de \textit{probes} e \textit{gauges} para monitorar o sistema e um mecanismo de adaptação para disparar mudanças com base nos dados coletados \cite{W1999392360}. Essa abordagem viria a influenciar diretamente a arquitetura MAPE-K, descrita na Seção \ref{sec:mapek}.

Em outubro de 2001, a IBM formalizou o conceito com seu manifesto de computação autônoma \cite{ibm2001}. A proposta central era que sistemas computacionais complexos deveriam possuir propriedades autônomas, ou seja, deveriam ser capazes de gerenciar independentemente as tarefas regulares de manutenção e otimização, reduzindo assim a carga de trabalho dos administradores. Além de cunhar o termo, a IBM sistematizou as quatro propriedades centrais de um sistema auto-gerenciável, discutidas na Seção \ref{sec:propriedades}.

Em 2003, a DARPA iniciou o programa \textbf{Self-Regenerative Systems (SPS)}, com foco em sistemas de computação militares capazes de manter funcionalidades críticas mesmo diante de erros não intencionais ou ataques. A abordagem incluía a geração de múltiplas versões funcionalmente equivalentes do software, tornando qualquer ataque incapaz de comprometer simultaneamente uma fração significativa das versões do programa. Um modelo de confiança era utilizado para afastar a computação de recursos propensos a causar dano, e uma arquitetura de replicação tolerante a intrusões de larga escala era desenvolvida em paralelo \cite{W1999392360}.

Por fim, em 2005, a NASA iniciou o projeto \textbf{Autonomous NanoTechnology Swarm (ANTS)}, cujo objetivo era explorar o cinturão de asteroides com um enxame de mil pequenas espaçonaves. Dado que entre 60\% e 70\% das naves eram esperadas a ser perdidas ao entrar no cinturão, as sobreviventes deveriam colaborar autonomamente para completar a missão. O projeto foi motivado também pelo longo atraso de comunicação entre sondas em espaço profundo e o controle em Terra, tornando o gerenciamento manual impraticável. A NASA já havia utilizado comportamentos autônomos nas missões DS1 (\textit{Deep Space 1}) e Mars Pathfinder, e o projeto ANTS representou um passo além nessa direção \cite{W1999392360}.

A Tabela \ref{tab:cronologia} apresenta uma síntese cronológica desses projetos e suas contribuições para o campo.

\begin{table}[H]
\caption{Cronologia de projetos influentes em auto-gerenciamento}
\label{tab:cronologia}
\centering
\begin{tabular}{lclp{7cm}}
\toprule
\textbf{Projeto} & \textbf{Ano} & \textbf{Origem} & \textbf{Contribuição principal} \\
\midrule
SAS   & 1997 & DARPA & Roteamento ad-hoc descentralizado e adaptativo \\
DASADA & 2000 & DARPA & Probes e gauges para monitoramento arquitetural \\
Computação Autônoma & 2001 & IBM & Formalização do termo e das quatro propriedades Self-X \\
SPS   & 2003 & DARPA & Sistemas auto-regenerativos resistentes a ataques \\
ANTS  & 2005 & NASA  & Enxame autônomo para exploração espacial \\
\bottomrule
\end{tabular}
\end{table}

\section{Propriedades de Auto-Gerenciamento}
\label{sec:propriedades}

As quatro propriedades fundamentais de um sistema autônomo, conforme sistematizadas pela IBM, são conhecidas como propriedades \textit{Self-X}. Cada uma aborda uma dimensão distinta do auto-gerenciamento \cite{kephart2003, W1999392360}.

\textbf{Auto-configuração (\textit{Self-Configuration}):} Um sistema autônomo configura-se de acordo com objetivos de alto nível. Isso significa que o administrador especifica \textit{o que} é desejado, não \textit{como} alcançá-lo. O sistema instala-se e ajusta-se automaticamente de acordo com as necessidades da plataforma e dos usuários \cite{W1999392360}.

\textbf{Auto-otimização (\textit{Self-Optimization}):} O sistema otimiza proativamente o uso de seus recursos, podendo decidir por iniciativa própria realizar mudanças para melhorar o desempenho ou a qualidade do serviço. Essa propriedade distingue-se pelo seu caráter \textit{pró-ativo}: o sistema não aguarda uma solicitação externa para agir, ele antecipa a necessidade \cite{W1999392360}.

\textbf{Auto-cura (\textit{Self-Healing}):} O sistema detecta e diagnostica problemas, que podem variar de erros de baixo nível em chips de memória a falhas de software em serviços de diretório. Sempre que possível, o sistema tenta corrigir o problema --- por exemplo, alternando para um componente redundante ou instalando atualizações de software --- garantindo que o processo de correção não introduza novos problemas \cite{W1999392360}. A tolerância a falhas é um aspecto central: o sistema deve ser \textit{reativo} a falhas ou a seus primeiros sinais.

\textbf{Auto-proteção (\textit{Self-Protection}):} O sistema protege-se tanto de ataques maliciosos quanto de usuários que inadvertidamente realizam alterações prejudiciais. O sistema ajusta-se autonomamente para alcançar segurança, privacidade e proteção de dados \cite{W1999392360}. Além de reagir a ameaças, o sistema deve ser capaz de antecipar brechas de segurança e impedir que ocorram, exibindo características \textit{proativas}.

Essas propriedades guardam estreita relação com as características de agentes de software identificadas por Wooldridge e Jennings, que incluem: (i) \textit{autonomia}, operar sem intervenção humana direta mantendo controle sobre suas ações e estado interno; (ii) \textit{habilidade social}, interagir com outros agentes por meio de algum tipo de linguagem de comunicação; (iii) \textit{reatividade}, perceber o ambiente e responder a mudanças em tempo hábil; e (iv) \textit{pró-atividade}, tomar iniciativa para atingir objetivos de forma dirigida por metas \cite{W1999392360}.

\section{O Laço de Controle Autônomo MAPE-K}
\label{sec:mapek}

Para operacionalizar as propriedades Self-X, a IBM propôs um modelo de referência para laços de controle autônomos denominado MAPE-K --- sigla para \textit{Monitor, Analyse, Plan, Execute over Knowledge} \cite{W1999392360, lalanda2013}. Este modelo descreve o ciclo funcional de um \textit{Elemento Autônomo}, composto por um \textit{Gerenciador Autônomo} acoplado a um \textit{Elemento Gerenciado}.

O \textit{Elemento Gerenciado} pode ser qualquer recurso de software ou hardware ao qual se deseja conferir comportamento autônomo: um servidor web, um banco de dados, um sistema operacional, um cluster de máquinas, uma rede com ou sem fio, uma CPU, etc. \cite{W1999392360}. A interação entre o gerenciador e o elemento gerenciado ocorre por meio de dois tipos de componentes de interface:

\begin{itemize}
    \item \textbf{Sensores (\textit{Sensors}/\textit{Probes}/\textit{Gauges}):} Coletam informações sobre o elemento gerenciado. Para um servidor web, isso pode incluir tempo de resposta a requisições de clientes, uso de rede, disco, CPU e memória \cite{W1999392360}. Os \textit{gauges} podem filtrar, agregar e processar os dados brutos coletados pelos \textit{probes} antes de reportá-los ao gerenciador, e podem residir em máquinas distintas do elemento gerenciado.

    \item \textbf{Efetores (\textit{Effectors}):} Implementam mudanças no elemento gerenciado. Mudanças podem ser de granularidade grossa --- como adicionar ou remover servidores de um cluster de servidores web --- ou de granularidade fina --- como alterar parâmetros de configuração individuais em um servidor \cite{W1999392360}.
\end{itemize}

\subsection{Monitoramento}
\label{subsec:monitoramento}

O monitoramento envolve a captura de propriedades do ambiente que são relevantes para as propriedades Self-X do sistema. São identificados dois tipos de monitoramento em sistemas autônomos \cite{W1999392360}:

\textbf{Monitoramento passivo:} Realizado sem modificar o sistema monitorado. No Linux, por exemplo, comandos como \texttt{top} retornam informações sobre a utilização de CPU por cada processo, enquanto \texttt{vmstat} fornece estatísticas de memória e CPU. O diretório \texttt{/proc} é um sistema de arquivos virtual que contém informações de tempo de execução como memória do sistema, informações de CPU, configuração de hardware, entre outros. Ferramentas similares existem para outros sistemas operacionais.

\textbf{Monitoramento ativo:} Envolve a instrumentação do software --- a modificação ou adição de código à implementação da aplicação ou do sistema operacional para capturar chamadas de função ou de sistema. Ferramentas especializadas como o ProbeMeister são capazes de inserir probes diretamente no bytecode Java compilado, automatizando esse processo.

\subsection{Planejamento e Adaptação}
\label{subsec:planejamento}

O Gerenciador Autônomo utiliza os dados coletados pelos sensores e o conhecimento interno do sistema para analisar o estado do elemento gerenciado, planejar as ações necessárias e executá-las por meio dos efetores \cite{W1999392360}. Os objetivos do gerenciador são normalmente expressos por meio de políticas, que podem assumir três formas principais:

\textbf{Políticas ECA (\textit{Event-Condition-Action}):} Definem ações diretas a serem tomadas quando um evento ocorre e uma condição é satisfeita. Por exemplo: ``quando 95\% dos servidores web apresentar tempo de resposta superior a 2 segundos e houver recursos disponíveis, aumentar o número de servidores ativos'' \cite{W1999392360}. Sua simplicidade é sua maior vantagem. Contudo, quando múltiplas políticas são especificadas, podem surgir conflitos entre regras que são difíceis de detectar antecipadamente e que só se manifestam em tempo de execução.

\textbf{Políticas de objetivo (\textit{Goal Policies}):} Especificam critérios que caracterizam estados desejáveis do sistema, deixando ao próprio sistema a tarefa de determinar como atingi-los. São mais expressivas que as políticas ECA, mas exigem maior capacidade de raciocínio e planejamento por parte do gerenciador \cite{W1999392360}.

\textbf{Funções de utilidade (\textit{Utility Functions}):} Definem quantitativamente o grau de desejabilidade de cada estado do sistema, atribuindo um valor numérico a cada combinação de parâmetros observáveis. Permitem ao gerenciador tomar decisões otimizadas mesmo quando o estado ideal não pode ser alcançado, identificando o estado menos indesejável entre as opções disponíveis \cite{W1999392360}.

Uma abordagem mais sofisticada de planejamento é o \textit{modelo arquitetural dirigido}, no qual o gerenciador mantém um modelo da arquitetura do sistema gerenciado. Este modelo é atualizado continuamente pelos dados dos sensores e utilizado para planejar adaptações. A vantagem central dessa abordagem é a possibilidade de verificar, antes da execução, se a mudança planejada preserva a integridade do sistema \cite{W1999392360}.

\subsection{Base de Conhecimento}
\label{subsec:conhecimento}

O componente de \textit{Conhecimento} fundamenta todas as fases do laço MAPE-K. Ele pode ser tão simples quanto um conjunto de regras estáticas definidas por especialistas humanos, ou tão sofisticado quanto modelos preditivos construídos a partir de técnicas de aprendizado de máquina \cite{W1999392360, lalanda2013}. Entre as abordagens mais relevantes, destacam-se:

\textbf{Aprendizado por reforço:} Aprende políticas de gerenciamento a partir da observação de ações em diferentes estados do sistema, avaliando as consequências de cada ação. Sua principal vantagem é não exigir um modelo explícito do sistema gerenciado, sendo adequado para ambientes altamente dinâmicos. A escalabilidade em espaços de estados grandes, no entanto, é um desafio reconhecido \cite{W1999392360}.

\textbf{Redes Bayesianas:} Utilizadas para a seleção probabilística de algoritmos e para a classificação sensível a custo de estados do sistema, auxiliando o gerenciador na tomada de decisões sob incerteza \cite{W1999392360}.

\textbf{Modelos arquiteturais:} Representações estruturais do sistema gerenciado, usadas para planejar adaptações e verificar a validade das mudanças antes de sua execução no sistema real, conforme descrito na seção anterior \cite{W1999392360, lalanda2013}.

A distinção entre o componente de planejamento e o de conhecimento é frequentemente tênue na prática: o conhecimento acumulado sobre o estado e o histórico do sistema gerenciado é o que torna possível ao gerenciador construir planos de adaptação confiáveis \cite{W1999392360}.
